\documentclass{article}
\usepackage{amsmath,amssymb}
\usepackage[vocab={numbers}]{leanify}
\newtheorem{lemma}{Lemma}
\newtheorem{definition}{Definition}
\begin{document}

\section{Constants with dependence}

Here and below, $K$ denotes a positive constant that depends only on $n$.

In the lemmas below, $A_k>0$ and $b_k\ge0$ may depend on $k$ and the fixed data, and are uniform in $N$ and $x$.

\begin{lemma}\label{powers}
For every positive integer $n$ and every $x\in[0,1]$, $x^n\le K$.
\end{lemma}

\begin{lemma}\label{exponential}
For every positive integer $k$ and every positive integer $N$, $N^k\le A_k2^N$.
\end{lemma}

\begin{lemma}\label{logarithm}
For every positive integer $k$, every positive integer $N$ and every $x\ge1$, $\log x\le A_kx^{1/k}$.
\end{lemma}

\begin{lemma}\label{sine}
For every positive integer $k$ and every $x\in\mathbb R$, $|\sin(kx)|\le b_k|x|$.
\end{lemma}

\begin{lemma}\label{noparameter}
For every positive integer $N$ and every $x\in[0,1]$, $x^N\le K$.
\end{lemma}

The constants $c$ and $c_m$ are positive and independent of $N$.

\begin{lemma}\label{independent}
For every positive integer $m$ and every positive integer $N$, $c_m\le N+1$.
\end{lemma}

A dependence that the article does not define is not read:

In the lemmas below, $D>0$ may depend on the shape of the domain.

\section{Multi-indices}

Define $f:\mathbb R^2\to\mathbb R$ by $f(x,y)\coloneqq e^{x}\sin y$.

In the lemmas below, $C_I>0$ may depend on $I$, and is uniform in $x$ and $y$.

\begin{lemma}\label{derivatives}
For every $I\in\mathbb N_0^2$ and every $x,y\in[0,1]$, $|\partial^If(x,y)|\le C_I$.
\end{lemma}

\begin{lemma}\label{explicit}
$\forall I\in\mathbb N_0^2\ \exists C_I: |\partial^If(0,0)|\le C_I$.
\end{lemma}

\begin{lemma}\label{components}
For every $I\in\mathbb N_0^2$, $I_1+I_2\ge0$.
\end{lemma}

\begin{lemma}\label{sums}
$\forall x\in\mathbb R\ \exists y,z: y+z=x$.
\end{lemma}

\section{Classes of functions}

For a positive integer $j$, we set $\lambda=2^{j}$ and $\mu=\lambda^{1/2}$.

In the definitions below, $B_I>0$ may depend on $I$, and is uniform in $j$.

\begin{definition}\label{bernstein}
A map $f:\mathbb R^2\to\mathbb R$ belongs to $\mathcal B_s$ if $|\partial^If(x,y)|\le B_I\lambda^{s+I_1+I_2}$ for every $I\in\mathbb N_0^2$ and every $x,y\in\mathbb R$.
\end{definition}

\begin{definition}\label{symbols}
A map $a:\mathbb R^2\to\mathbb R$ belongs to $\mathcal S_m$ if $|\partial^Ia(x,\xi)|\le B_I(1+|\xi|)^{m-I_2}$ for every $I\in\mathbb N_0^2$ and every $x,\xi\in\mathbb R$.
\end{definition}

\begin{lemma}\label{bernstein-inclusion}
For all $s,t\in\mathbb R$ with $s\le t$ and every $f\in\mathcal B_s$, $f\in\mathcal B_t$.
\end{lemma}

\begin{lemma}\label{symbol-inclusion}
For all $m,p\in\mathbb R$ with $m\le p$ and every $a\in\mathcal S_m$, $a\in\mathcal S_p$.
\end{lemma}

\end{document}
